\paragraph{Why HG002.} HG002 (NA24385) is the son in the Ashkenazi Jewish
trio of the Genome in a Bottle consortium, a Personal Genome Project
participant with open consent and the most thoroughly characterised human
genome~\cite{zook2016,wagner2022}. Three properties make it a good test bed
for an annotation paper: a curated benchmark call set
(\NVariants{} small variants on chromosomes 1--22 after decomposition) with
explicit benchmark regions; matching benchmark VCFs for both parents
(HG003, HG004), which allow inheritance to be checked; and a founder
population, in which global allele frequencies can be misleading. The
question we ask is the classic one of MacArthur et
al.~\cite{macarthur2012}: \emph{how many genes does a healthy person knock
out, which of these calls are real, and what does this genome carry that a
clinician would care about?} Each step of the answer needs a different kind
of annotation, which is the point of the vepyr design.

\paragraph{1. A healthy genome carries tens of real LoF variants, and many
more false ones.} The core annotation finds \NPlof{} variants with a pLoF
consequence (stop-gained, frameshift, canonical splice site) in a MANE
Select transcript~\cite{morales2022}. Filters in the spirit of
LOFTEE~\cite{karczewski2020} reduce them to \NHc{} high-confidence pLoF
variants in \NHcGenes{} genes, \NHcHom{} of them homozygous, in line with
earlier estimates of about a hundred LoF variants and about twenty
homozygous LoF genes per genome~\cite{macarthur2012}, given that our counts
are restricted to MANE transcripts and benchmark regions. The removed calls
form recognisable classes, and each class comes from a different tier:
\begin{itemize}
\item \emph{Tier 1 (core cache).} \FlagLastExon{} premature stops in the
  last exon, which escape nonsense-mediated decay. \FlagRefMinor{} calls
  where the ``LoF'' allele is near-fixed in gnomAD: pairs of nearby
  frameshifts in \textit{SON} and \textit{SLAIN1} that restore the reading
  frame (GRCh38 carries a rare sequence there), and \textit{FMO2}
  p.Gln472Ter, the \textit{FMO2*2} allele that truncates FMO2 in most
  humans~\cite{dolphin1998}---a real, but species-wide, loss.
\item \emph{Tier 2 (plugins).} \FlagClinvarBenign{} variants classified
  benign in ClinVar, including \textit{LPL} p.Ser474Ter (S447X), a stop-gain
  associated with lower triglycerides~\cite{rip2006}, and the
  \textit{CYP2D6*4} splice-acceptor allele, a genuine loss of function in a
  pharmacogene that ClinVar classifies as a drug-response allele.
  \FlagSpliceai{} canonical splice-site variants without SpliceAI
  support~\cite{jaganathan2019}.
\item \emph{Tier 3 (ad hoc).} \FlagSegdup{} calls in segmental duplications
  (olfactory receptors, \textit{NBPF}, \textit{GOLGA8}, \textit{NPIPB},
  \textit{CYP2D6}) and \FlagLowmap{} in low-mappability regions, from GIAB
  stratifications~\cite{dwarshuis2024}; \FlagFiftybp{} premature stops within
  50\,bp of the last exon junction, from Ensembl exon coordinates.
\end{itemize}

\paragraph{2. Purifying selection is visible in one genome.} Joining the
frame with gnomAD v4.1 gene constraint (read from the gnomAD bucket by URL)
shows that only \LoeufAllPlofPct\% of HG002 pLoF variants fall in the 20\%
most LoF-intolerant genes, against \LoeufSynPct\% of rare synonymous variants
($p=\LoeufAllPlofP$), and none of the rare high-confidence ones does. A
healthy individual carries LoF variants, but in genes that tolerate them.
The rare synonymous reference matters: constrained genes are longer, so a
naive uniform expectation (20\%) would understate the depletion.

\paragraph{3. One clinically classified variant: an Ashkenazi founder
allele.} The only ClinVar pathogenic or likely pathogenic variant with
review criteria is \textit{F11} \PlpHgvsp{}, the factor~XI ``type~II''
allele (Glu117Ter in legacy numbering), one of the two founder mutations
that account for most factor~XI deficiency in Ashkenazi
Jews~\cite{asakai1991}. HG002 is heterozygous; the allele was transmitted
by the \PlpParent{}; its gnomAD frequency is \PlpAf\% overall but \PlpAsj\%
in Ashkenazi Jewish samples (\PlpFold-fold). The same founder effect shows
across the rare pLoF variants: \NRareAsjCommon{} of the \NHcRare{} that are
rare in gnomAD overall exceed 1\% in Ashkenazi Jewish samples, so
``rare'' depends on the reference population. Population-specific
frequencies come from the core cache; ClinVar from a plugin cache; the
parental genotypes from the parents' VCFs streamed from S3.

\paragraph{4. The trio resolves origins and exposes one false homozygous
LoF.} Nine of the ten rare high-confidence pLoF variants have a resolved
origin (\NRareMaternal{} maternal, \NRarePaternal{} paternal, one carried by
both parents) and none is a candidate \emph{de novo} variant. The one
homozygous rare pLoF, a 45-bp deletion of the canonical exon~2 donor of
\textit{PEX5}, would imply Zellweger spectrum disorder if it were a loss of
function. Both parents are homozygous for it, which already argues against
rarity (ClinVar classifications of the deletion are conflicting), and the exon has an alternative donor 45\,bp downstream that keeps
the reading frame. Absence from gnomAD here reflects the length of the
deletion, not its frequency. The case is a useful reminder that
transcript-level LoF calls need biological context; in this analysis the
context comes from the ad hoc tier.

\paragraph{5. What this genome cannot tell us.} Interval coverage of MANE
coding sequence by the benchmark regions shows that \NClinGenLtNinety{}
ClinGen disease genes~\cite{strande2017} are less than 90\% assessable in
HG002 v4.2.1, \NClinGenLtFifty{} less than half. They include genes in
segmental duplications and falsely duplicated regions of GRCh38
(\ClinGenFalseDup{}; \textit{CBS}, classic homocystinuria, is a known
example~\cite{aganezov2022}), \textit{FXN} (Friedreich ataxia) and GC-rich
genes such as \textit{ETFB}, \textit{GNPTG} and \textit{FAM20C}. Finding no
pathogenic variant in these genes is not evidence of absence.

\paragraph{Why this belongs in the vepyr paper.} The analysis is a single
Python session: \AnnotateMin\,min of annotation of the whole genome with
VEP-identical core fields and four plugins, then \AnalyzeMin\,min for
everything else, including \RemoteSec\,s to read \NRemoteSources{} remote
sources (\RemoteRows{} million records) without downloading, indexing or
converting them. Three kinds of annotation with three update rates---Ensembl
releases, ClinVar weeks, per-question data---are handled by three
mechanisms: a core cache, plugin caches and DataFrame operations on data read
from URLs. The biology is not new, which is the point: a known result
(depletion of LoF from constrained genes), a known founder allele and a
known reference artefact are recovered quickly, reproducibly and in one
place, which is what users need when the question, not the method, is new.
